\subsection{Формула Тейлора для $e^x, \sin x, \cos x$. Оценка остатка.}

Формула Тейлора для функции $f(x)$ в окрестности произвольной точки $a$:
\[ f(x) = \sum_{k=0}^n \frac{f^{(k)}(a)}{k!}(x-a)^k + R_n(x) \]
Остаточный член в форме Лагранжа:
\[ R_n(x) = \frac{f^{(n+1)}(\xi)}{(n+1)!}(x-a)^{n+1}, \quad \xi \in (a, x) \]

Обычно для этих функций используют разложение в точке $a=0$ (формула Маклорена), так как оно имеет наиболее простой вид.

\subsubsection*{1. Функция $f(x) = e^x$}

\textbf{Общий случай (точка $a$):}
\begin{enumerate}
	\item $n$-я производная: $f^{(n)}(x) = e^x$.
	\item Значение в точке $a$: $f^{(n)}(a) = e^a$.
	\item Формула:
    \[ e^x = e^a + e^a(x-a) + \frac{e^a}{2!}(x-a)^2 + \dots + \frac{e^a}{n!}(x-a)^n + R_n \]
    \[ e^x = e^a \sum_{k=0}^n \frac{(x-a)^k}{k!} + \frac{e^\xi}{(n+1)!}(x-a)^{n+1} \]
\end{enumerate}

\textbf{Частный случай ($a=0$):}
\[ e^x = 1 + x + \frac{x^2}{2!} + \dots + \frac{x^n}{n!} + \frac{e^\xi}{(n+1)!}x^{n+1} \]

\textbf{Оценка остатка:}
Для любого фиксированного $x$ и фиксированного $a$: $|e^\xi|$ ограничено константой $M$ на отрезке между $a$ и $x$.
\[ |R_n| \le M \frac{|x-a|^{n+1}}{(n+1)!} \xrightarrow{n \to \infty} 0 \]

\subsubsection*{2. Функция $f(x) = \sin x$}

1.  \textbf{Производная $n$-го порядка:}
    \[ f^{(n)}(x) = \sin\left(x + \frac{\pi n}{2}\right) \]
2.  \textbf{Формула Тейлора в точке $a$:}
    \[ \sin x = \sum_{k=0}^n \frac{\sin(a + \frac{\pi k}{2})}{k!} (x-a)^k + R_n(x) \]
    
    Это выражение громоздкое, поэтому обычно приводят \textbf{разложение в нуле ($a=0$)}:
    Значения производных в нуле: $0, 1, 0, -1, \dots$
    \[ \sin x = x - \frac{x^3}{3!} + \frac{x^5}{5!} - \dots + (-1)^m \frac{x^{2m+1}}{(2m+1)!} + R_{2m+1} \]

\textbf{Оценка остатка:}
Так как $|\sin(\dots)| \le 1$ и $|\cos(\dots)| \le 1$:
\[ |R_n(x)| = \left| \frac{f^{(n+1)}(\xi)}{(n+1)!}(x-a)^{n+1} \right| \le \frac{|x-a|^{n+1}}{(n+1)!} \]
Предел этого выражения при $n \to \infty$ равен 0 для любого $x$.

\subsubsection*{3. Функция $f(x) = \cos x$}

1.  \textbf{Производная $n$-го порядка:}
    \[ f^{(n)}(x) = \cos\left(x + \frac{\pi n}{2}\right) \]
2.  \textbf{Формула Тейлора в точке $a$:}
    \[ \cos x = \sum_{k=0}^n \frac{\cos(a + \frac{\pi k}{2})}{k!} (x-a)^k + R_n(x) \]

    \textbf{Разложение в нуле ($a=0$):}
    Значения производных в нуле: $1, 0, -1, 0, \dots$
    \[ \cos x = 1 - \frac{x^2}{2!} + \frac{x^4}{4!} - \dots + (-1)^m \frac{x^{2m}}{(2m)!} + R_{2m} \]

\textbf{Оценка остатка:}
Аналогично синусу:
\[ |R_n(x)| \le \frac{|x-a|^{n+1}}{(n+1)!} \xrightarrow{n \to \infty} 0 \]
